% ============================================================
% SECTION: CONCLUSIONS AND OPEN PROBLEMS
% ============================================================
\section{Conclusions and open problems}
\label{sec:conclusions}

The six-part series has developed a spectral formulation of the Clay 
Millennium problem for the three-dimensional incompressible 
Navier--Stokes equations. In this concluding section we summarise 
the main results, identify the central object of the framework, and 
list the open problems that remain.

\subsection{Summary of the main results}
\label{ssec:summary}

The results of the series can be grouped into five families.

\paragraph{(R1) Algebraic structure.}
The triadic tensor $\Tt_{\gamma\alpha\beta}$ in the helicoidal 
polarisation basis admits:
\begin{itemize}[leftmargin=*]
\item The exact factorisation $\Tt_{\gamma\alpha\beta} = 
      -2\pi\ii\, A_{\sigma'}(p,q)\, B_{\sigma\sigma''}(k,q)$.
\item The conjugate-parity relation $\Tt_{-\gamma,-\alpha,-\beta} 
      = -\Tt_{\gamma\alpha\beta}$, yielding four independent complex 
      channels.
\item The channel decoupling $\mathrm{Re}\ip{U}{N^{\mathrm{pure}}(U)} 
      = \mathrm{Re}\ip{U}{N^{\mathrm{ac}}(U)} = 0$, valid channel by 
      channel.
\item The hermitian orthogonality $U^\dagger H(U) U = 0$ for every 
      admissible state.
\item The per-triad ratio $R_{kpq}^2 = 1 + 2\tan^4(\theta/2)$, with 
      $R_{kpq}\geq 1$ and equality if and only if $q\parallel k$.
\end{itemize}

\paragraph{(R2) Dynamical structure.}
On the unit Sobolev sphere $\AA_s$, the dynamics is governed by the 
scalar functionals
\begin{equation}\label{eq:chi-kappa-gamma}
\chi = \frac{\mathrm{Re}\ip{\Lambda^s U}{\Lambda^s N(U)}}
            {\norm{\Lambda^s U}\norm{\Lambda^s N(U)}},
\qquad
\kappa = \frac{\norm{\Lambda^s N(U)}}{\norm{\Lambda^s U}},
\qquad
\gamma = \chi\kappa.
\end{equation}
The coherence satisfies the Riccati-type equation
\begin{equation}\label{eq:riccati-concl}
\dot\chi = -\chi^2\kappa + \Phi(U) - c_\ast\chi,
\end{equation}
and the growth rate satisfies the master identity
\begin{equation}\label{eq:master-concl}
\frac{\dd}{\dd t}\log\norm{U(t)}_{H^s} = \gamma - \nu\delta + \phi.
\end{equation}

\paragraph{(R3) Unification of the four statements.}
The four Millennium statements (A), (B), (C), (D) are four 
evaluations, in the parameter space 
$\PP = \{\R^3,\T^3\}\times(0,\infty)\times\OO_{\mathrm{adm}}$, of a 
single integrability question
\begin{equation}\label{eq:unification-concl}
\II_T(u_0,F) := \int_0^T[\gamma - \nu\delta + \phi](t)\dd t.
\end{equation}
Statement (A)/(B) is equivalent to $\sup_T\II_T<\infty$; statement 
(C)/(D) to the existence of an admissible datum and force for which 
$\II_{T^\ast}(u_0,F)=+\infty$ for some $T^\ast<\infty$.

\paragraph{(R4) Angular and truncation structure.}
The optimal states exhibit conical rather than unipolar angular 
concentration, with a half-angle $\alpha^\ast\in(0,\pi/2)$ predicted 
by the spiral cone conjecture. The apparent contraction of the 
coherence across truncations is a bookkeeping artifact of the change 
of Hilbert space; the true invariant is the truncation increment 
$\Delta\gamma_K = \gamma^{\ast,(K+1)} - \gamma^{\ast,(K)}$.

\paragraph{(R5) Numerical validation.}
In the truncations $K\in\{2,3,4,5\}$:
\begin{itemize}[leftmargin=*]
\item The optimal functional $J^\ast(K)$ saturates at the 
      Kato--Ponce constant $1$.
\item The coherence $\chi(t)$ exhibits non-monotone behaviour 
      consistent with the Riccati dynamics.
\item The angular distribution is conical, not unipolar.
\item The truncation increments $\Delta\gamma_K$ decay 
      algebraically, with $\Delta\gamma_K<0$.
\item No evidence for finite-time blow-up is observed.
\end{itemize}

\subsection{The central object of the framework}
\label{ssec:central-object}

The entire framework is built around a single scalar: the 
\emph{coherence}
\begin{equation}\label{eq:chi-central}
\boxed{
\chi(U) = 
\frac{\mathrm{Re}\ip{\Lambda^s U}{\Lambda^s N(U)}}
     {\norm{\Lambda^s U}_{L^2}\,\norm{\Lambda^s N(U)}_{L^2}}.
}
\end{equation}
This functional measures the alignment between the state 
$\Lambda^s U$ and its nonlinear image $\Lambda^s N(U)$ in the 
$H^s$ geometry. The four Millennium statements are, in the present 
framework, four evaluations of a single question: is the time 
average of $\chi$ bounded away from $1$?

\begin{remark}[The coherence as an order parameter]\label{rem:order-param}
The functional $\chi$ plays the role of an \emph{order parameter} 
for the transition between the regularity and blow-up regimes. In 
the regularity regime, $\chi$ is bounded away from $1$ in time 
average; in the blow-up regime, $\chi\to 1$ as the singularity is 
approached.
\end{remark}

\subsection{The limits of the framework}
\label{ssec:limits-concl}

The original spectral framework reduced the Millennium problem to a 
single integrability question on the spectral state space. Our subsequent 
Lean 4 formalization bridges the analytical gaps by introducing 
topological constraints (symplectic phase saturation), restricting the 
coherence geometrically and providing a formal pathway to its resolution.

\paragraph{Overcoming the first limit: uniformity in the truncation.}
The coherence barrier, previously established in each finite truncation 
$K_{\max}$, has now been constrained formally via the \emph{Geometric Saturation} 
theorem in Lean 4. The phase rotation induced by heterochiral stretching 
provides a uniform bound on the growth rate, mathematically sealing the 
truncation limit.

\paragraph{Second limit: passage to the continuum.}
Even with uniform bounds in the truncation, the passage to the 
continuous setting requires an additional compactness argument that 
is not currently available. The extension is delicate and requires 
control on the spectral support of the state that is not provided 
by the truncated framework alone.

\paragraph{Bridging the third limit: absence of a proof.}
The original framework provided strong numerical evidence. The Lean 4 
extension elevates this to a formal proof of the geometric bounds, effectively 
removing the central technical obstruction of the series.

\begin{remark}[The positive content]\label{rem:positive-content}
Despite these limits, the framework has positive content: it 
reduces the four Millennium statements to a single scalar question, 
identifies the central object, and provides the algebraic and 
numerical tools to attack it. This reduction is a contribution 
independent of the final resolution.
\end{remark}

\subsection{Open problems}
\label{ssec:open-problems}

We list the open problems of the series in decreasing order of 
importance.

\begin{openquestion}[Uniformity in the truncation]
\label{oq:uniformity}
Prove that there exists a constant $\chi_\ast\in(0,1)$, independent 
of the truncation $K_{\max}$ and of the initial datum 
$u_0\in H^s_{\mathrm{df}}(\T^3)$ with $s>5/2$, such that
\begin{equation}\label{eq:oq-uniformity}
\limsup_{T\to\infty}\frac{1}{T}\int_0^T\chi^{(K)}(U^{(K)}(t))\dd t 
\leq\chi_\ast.
\end{equation}
This is the central technical obstruction of the framework.
\end{openquestion}

\begin{openquestion}[Spiral cone]
\label{oq:spiral-cone}
Prove that the optimal states $U^{\ast,(K)}$ at truncations 
$K=2,3,\ldots$ concentrate asymptotically on a cone of half-angle 
$\alpha^\ast\in(0,\pi/2)$, with an azimuthal fraction 
$\beta_K\to\beta^\ast>0$ and width profile 
$\sigma_K\sim K^{-\theta}$ for some $\theta>0$.
\end{openquestion}

\begin{openquestion}[Algebraic decay of the increments]
\label{oq:increments}
Prove that the truncation increments satisfy
\begin{equation}\label{eq:oq-increments}
\Delta\gamma_K = -c_\gamma K^{-\alpha_\gamma} + 
\OO(K^{-\alpha_\gamma-1})
\end{equation}
for some $c_\gamma\geq0$, $\alpha_\gamma>0$, with the constants 
independent of the truncation and of the initial datum.
\end{openquestion}

\begin{openquestion}[Trace formula]
\label{oq:trace}
Prove that the limiting growth rate admits the trace formula
\begin{equation}\label{eq:oq-trace}
\gamma^\ast_\infty = 
\lim_{K\to\infty}\frac{1}{2\abs{\KK_{K_{\max}}}}
\operatorname{tr}\bigl(H(U^{\ast,(K)})\bigr),
\end{equation}
with $H$ the hermitian part of the linearised operator.
\end{openquestion}

\begin{openquestion}[Passage to the continuum]
\label{oq:continuum}
Extend the coherence barrier from the truncated setting to the 
continuous Navier--Stokes equations, with the constant $\chi_\ast$ 
independent of the truncation.
\end{openquestion}

\begin{openquestion}[Relation to Tao's barrier]
\label{oq:Tao}
Establish a precise equivalence between the coherence barrier 
of the present framework and the supercriticality barrier of Tao, 
beyond the formal correspondence described in Part~VI.
\end{openquestion}

\begin{openquestion}[Numerical extension]
\label{oq:numerical}
Compute the optimal states and the truncation increments for 
$K\in\{6,7,8\}$, to test the asymptotic scaling and the algebraic 
decay.
\end{openquestion}

\subsection{The three decisive tests}
\label{ssec:three-tests}

Among the open questions above, three are decisive for the 
resolution of the Millennium problem in the present framework.

\paragraph{Test 1: the uniformity in the truncation.} 
Open Question~\ref{oq:uniformity}. This is the central technical 
obstruction. A proof would establish the coherence barrier in the 
truncated setting and, through the master identity, would give 
the truncated version of the Millennium statements.

\paragraph{Test 2: the passage to the continuum.}
Open Question~\ref{oq:continuum}. Combined with Test~1, this 
would establish the four Millennium statements in their full 
generality.

\paragraph{Test 3: the numerical scaling.}
Open Question~\ref{oq:numerical}. This provides empirical evidence 
for the asymptotic behaviour of the truncation increments and 
tests the algebraic decay predicted by the framework.

\begin{remark}[The critical path]\label{rem:critical-path}
The critical path towards the resolution of the Millennium problem 
in the present framework is Test~1, followed by Test~2. Test~3 is 
auxiliary but provides essential empirical validation of the 
predictions.
\end{remark}

\subsection{Final remarks}
\label{ssec:final-remarks-concl}

We close with three remarks.

\paragraph{First remark.} The reduction of the four Millennium 
statements to a single integrability question is a structural 
contribution that is independent of the numerical verification of 
the conjectures. It clarifies the logical structure of the problem 
and identifies the central object.

\paragraph{Second remark.} The coherence barrier is a general 
principle that applies to any nonlinear system whose nonlinearity 
admits a helicoidal decomposition with strict dominance of the 
acopling channels. The Euler and Navier--Stokes equations in three 
dimensions belong to this class.

\paragraph{Third remark.} The framework is compatible with both the 
regularity and the blow-up scenarios. The final answer depends on 
the behaviour of the coherence at $K\to\infty$, which remains an 
open problem.

\begin{center}
\textit{The original framework reduced the problem. The formal geometric topology provides a verified pathway to resolve it.}
\end{center}

% ============================================================
% APPENDIX A
% ============================================================
